% Generated draft: requires mathematical and source review.
\documentclass[11pt]{article}
\usepackage[margin=25mm]{geometry}
\usepackage{fontspec}
\setmainfont{Latin Modern Roman}
\usepackage{amsmath,amssymb,mathtools,unicode-math}
\setmathfont{Latin Modern Math}
\usepackage{enumitem,needspace,xurl,hyperref}
\setlength{\parindent}{0pt}
\setlength{\parskip}{5pt}
\setlength{\emergencystretch}{3em}
\newcommand{\N}{\mathbb{N}}
\newcommand{\Z}{\mathbb{Z}}
\newcommand{\Q}{\mathbb{Q}}
\newcommand{\R}{\mathbb{R}}
\newcommand{\C}{\mathbb{C}}
\newcommand{\F}{\mathbb{F}}
\newcommand{\A}{\mathbb{A}}
\newcommand{\PP}{\mathbb{P}}
\newcommand{\E}{\mathbb{E}}
\newcommand{\eps}{\varepsilon}
\newcommand{\dd}{\,\mathrm{d}}
\begin{document}
\section*{UnsolvedMath ID: 30005734}
\textbf{Source problem code:} catalogue entry
\subsection{Definitions and mathematical statement}
The supplied record asks for an extension of a theory, but does not specify the results to be extended. Consequently, it determines a research problem rather than a fully quantified conjecture.

For interpreting the displayed equations, write \(\mathbb T^d=\R^d/\Z^d\), with \(d\geq 1\), and interpret \(c\in\R^d\) as a constant covector. The customary autonomous Tonelli framework takes a Hamiltonian \(H\colon\mathbb T^d\times\R^d\to\R\) of class \(C^2\), with positive definite momentum Hessian \(D^2_{pp}H(x,p)\), and uniform superlinearity:
\[
\lim_{R\to\infty}\ \inf_{\substack{x\in\mathbb T^d\\ |p|\geq R}}
\frac{H(x,p)}{|p|}=+\infty.
\]
These are explicit conventional interpretations of the notation and the word Tonelli; the supplied excerpt does not state the precise regularity and convexity conventions intended by the report.

The mechanical equation appearing in the record is
\[
H_{\mathrm{mech}}(x,Du(x)+c)
=\frac12|Du(x)+c|^2+V(x)=\alpha[c],
\qquad x\in\mathbb T^d,
\]
and the proposed generalization is
\[
H(x,Du(x)+c)=\alpha[c],
\qquad x\in\mathbb T^d.
\]
Here \(u\) is an unknown function on the torus, \(V\) is the potential in the mechanical case, and \(\alpha[c]\) denotes the scalar associated with \(c\) in the source notation. The record does not define \(\alpha\), specify the admissible potentials, or state the solution concept and regularity required of \(u\). In particular, interpreting \(u\) as a viscosity solution requires confirmation from the underlying report. Working on the torus expresses periodicity; no additional boundary condition is supplied.

For a locally Lipschitz function \(u\), its singular set can conventionally be written as
\[
\operatorname{Sing}(u)=\{x\in\mathbb T^d: u\text{ is not differentiable at }x\}.
\]
For a measurable trajectory \(\gamma\colon[0,\infty)\to\mathbb T^d\), its time-averaged occupation probability measure is
\[
\mu_T=\frac1T\int_0^T\delta_{\gamma(t)}\,\mathrm dt,
\qquad T>0.
\]
A limiting occupation measure in this position-space convention is a probability measure \(\mu\) such that, for some sequence \(T_j\to\infty\),
\[
\int_{\mathbb T^d}f\,\mathrm d\mu_{T_j}
\longrightarrow
\int_{\mathbb T^d}f\,\mathrm d\mu
\quad\text{for every }f\in C(\mathbb T^d).
\]
This definition explains the terminology but does not identify the dynamics intended in the source. The record does not specify the trajectories, whether occupation measures instead live in phase space, or the definition of the critical set, particularly at singular points.

\textbf{Requested extension.} Extend the relationships between critical sets, singular sets, and limiting occupation measures known for the displayed mechanical equation on \(\mathbb T^d\) to the displayed equation for general Tonelli Hamiltonians. The exact relationships, their hypotheses, the relevant dynamics, and their quantifiers are absent from the supplied record. A precise theorem or conjecture cannot therefore be stated without additional source information; the conventions above do not resolve those omissions.
\subsection{Short English statement}
The problem asks to extend relationships among critical points, singularities of solutions, and long-time occupation measures from mechanical Hamilton--Jacobi equations on a torus to general Tonelli Hamiltonians. The supplied record omits the specific relationships and several necessary definitions, so this is a statement-only research-program draft requiring mathematical review.
\subsection{Sources}
\begin{itemize}
\item Supplied local UnsolvedMath problem registry, permanent ID \(30005734\), title: \emph{Occupational Measures for Tonelli Hamilton--Jacobi Equations}; exact source problem code: \texttt{catalogue entry}. The record belongs to the Oberwolfach Reports Open Problems collection. No individual proposer is identified; the supplied proposed year is 2023.
\item Supplied report citation: \emph{Control Methods in Hyperbolic PDEs} (2024), Oberwolfach Reports, pp.\ 2941--3011, DOI \texttt{10.4171/owr/2023/52}. \url{https://doi.org/10.4171/owr/2023/52}. This is also listed among the literature sources as Oberwolfach Report 52/2023. No location within the report is supplied for the target passage.
\item Supplied literature source: Cannarsa and Cheng, \emph{Generalized characteristics and Lax--Oleinik operators}. \url{https://doi.org/10.1007/s00526-017-1219-4}.
\item The supplied literature assessment was checked on \texttt{2026-08-22}. It labels the problem open, describes the general Tonelli extension as requiring new ideas, and states that existing generalized-characteristic and Lax--Oleinik results cover narrower settings. This is supplied dated curation derived from the OWR record, not an independently verified assessment. The record itself calls for independent verification of time-sensitive status and cited claims.
\item The registry labels the statement \texttt{corrected\_verified} and reports a check against the official Oberwolfach report, asserting preservation of its hypotheses, notation, and target. These are supplied verification claims, not verification performed for this draft. The field designated as the original statement contains only an incomplete sentence describing a challenging open question and asserting that report context was retained verbatim; it does not supply the necessary mathematical context.
\end{itemize}
\end{document}
